\section{Archived learning and candidate-reuse controls}
\label{app:legacy_controls}
These controls come from the archived experiments.
We reaggregate the saved per-trial records and inspect their source
protocols; no new training or QPU execution is claimed here. Three
CSVs and their hashes accompany the reproduction script. These results
are separate from the compiler benchmarks.

\subsection{Matched classical feature learning}
The supervised control uses five qubits, ten trainable PQK parameters
and five local Z expectations. The classical map has five features
$z_j(x)=\cos(\alpha_j r_j^Tx+\beta_j)$, with fixed unit-norm
projections $r_j$ and ten trainable scales and biases. Both maps use
the same Gaussian kernel with $\gamma=1$, kernel-target alignment,
50 finite-difference optimization steps, learning rate 0.3,
minibatch size 16, and a precomputed-kernel SVM.

The datasets are one fixed 80-point two-moons realization (noise 0.15,
generation seed 3), and Wine classes 0 and 1 (130 points, first five
features). For each of ten paired randomized trials, the 60/40
stratified split precedes scaling. Scaling is fit on training inputs
alone; test inputs are transformed and clipped to $[0,\pi]$.
Trials vary split, initialization and minibatch randomness; they are
not independent newly sampled datasets. The models share the training
budget and split protocol, not identical initial feature maps.

\begin{table}[tb]
\caption{Archived matched feature-learning control. Accuracy is mean
$\pm$ sample SD over ten trials, in percent; KTA is the mean before
and after training. These are descriptive small-benchmark results,
not evidence of quantum advantage or a statistical equivalence claim.}
\label{tab:legacy_classical}
\begin{tabular}{llcc}
\toprule
Data & Map & KTA before/after & Accuracy\\
Moons & Classical & .2892/.3974 & $90.63\pm6.75$\\
      & PQK       & .3076/.3335 & $91.25\pm5.47$\\
Wine  & Classical & .2212/.3497 & $91.35\pm3.65$\\
      & PQK       & .3175/.4076 & $93.46\pm2.60$\\
\bottomrule
\end{tabular}
\end{table}

\subsection{Matched embedding placement}
The second control changes only the order of the same gate multiset.
Let D denote a data block and P a parameter block. Chronologically,
the Type II circuit is D-D-P-P-D and the interleaved Type III circuit
is D-P-D-P-D. Both use five qubits, three identical data blocks,
two identical parameter blocks and 20 trainable parameters. Logical
block depth is matched before compilation; equal physical compiled
cost is not assumed. These are fidelity kernels, not the preceding
projected kernels.

The datasets and train-only scaling follow the preceding protocol.
The ten trials pair initial parameters, split and minibatch order.
Training uses 30 KTA steps, learning rate 0.2, batch size 16 and
central finite differences of step $10^{-4}$. Nonzero gradients occur
in both structures (Table~\ref{tab:legacy_order}); interleaving does
not produce a universal accuracy ordering. The results do not assert
that every Type II or Type III circuit is trainable, nor that a
state-QFI rank has been recovered.

\begin{table}[tb]
\caption{Archived embedding-placement control, ten paired trials.
Initial gradient norm and post-training KTA are means. Accuracy is
mean $\pm$ sample SD in percent. Equal means do not establish equal
predictions, statistical equivalence, or equality on other tasks.}
\label{tab:legacy_order}
\begin{tabular}{llccc}
\toprule
Data & Order & $\|\nabla_\theta\mathrm{KTA}\|$ & KTA & Accuracy\\
Moons & II  & .0589 & .3544 & $94.69\pm3.62$\\
      & III & .0561 & .3455 & $94.69\pm3.62$\\
Wine  & II  & .1253 & .3950 & $94.04\pm3.20$\\
      & III & .1086 & .3878 & $94.04\pm3.68$\\
\bottomrule
\end{tabular}
\end{table}

\subsection{Terminal-heavy fidelity candidate reuse}
An earlier fixed candidate set varies the input-independent output
block in $U_\theta(x)=W_\theta E(x)$ across 36 circuits. The fidelity
overlap cancels $W_\theta$, so an embedding-specific key permits reuse
within four conservative groups. Four keys do not prove four minimal
inequivalent kernels: further equivalences between embeddings are
not excluded. This is a static candidate-set audit, not a test of
adaptive search with later data gates.

The saved replay has 1,692 candidate--category records (36 candidates
in 47 categories), each with four support and twenty query inputs.
Evaluating one representative per key avoids 32 of 36 candidate
evaluations, or 88.9\%. The unbatched support--query entry schedule
changes from 135,360 to 15,040 pair tasks. The maximum recorded
kernel discrepancy is $6.00\times10^{-15}$, and all recorded
stabilized AUROCs agree within a group. The algebraic suffix identity
justifies equality; the finite replay only checks this implementation.
Multiplying those task counts by an assumed shot count would be a
workload model, not measured shot savings or equal-precision evidence.

The associated safe reuse procedure is to identify an exactly
certified common suffix, construct a key including the residual
embedding, preprocessing and kernel specification, and cache only
the current candidate's evaluations. Any subsequent data-dependent
extension must be re-certified. A state-QFI computation is neither
required for suffix cancellation nor a substitute for this
kernel-equivalence check. Because an ordinary suffix baseline already
finds this reduction, it is not incremental evidence for the new
internal certificate or a contradiction of the later zero-increment
PQK resource surveys.

\subsection{Hardware provenance and mitigation scope}
The archived hardware CSV has 141 rows (47 categories at 2, 4 and
8 qubits) and reports a mean eight-qubit support--query correlation
near 0.125. This is not a mitigated result. In the inspected export,
the backend field is redacted and matching raw counts, calibration
records and job identifiers have not been established. The inspected
legacy script constructs compute--uncompute fidelity circuits for
support--support entries; it does not by itself establish the
provenance of all support--query values in the CSV. We therefore do
not use that summary as a verified hardware resource endpoint here.

Measurement mitigation such as M3~\cite{m3} or Twirled Readout Error
eXtinction (TREX)~\cite{ibmtrex} could be assessed in a separately
specified experiment. TREX targets expectation-value readout errors
through measurement twirling and calibration; applying it to a
projected-feature experiment is not the same experiment as correcting
compute--uncompute counts. Neither method establishes recovery from
all gate and decoherence errors. Calibration effort and estimator
variance must be included, and improvement in kernel accuracy or
downstream classification must be measured rather than assumed.
